\section{Chapitre~\ref{sec:code}. Le code Morse}

Les bornes du chapitre relèvent de la théorie de l'information et du calcul; sont mesurés l'endroit où naît le bruit dans le canal, la vitesse du cerveau et le coût d'une impulsion, tandis que la question du code que le cortex utilise réellement reste ouverte.

{\footnotesize
\setlength{\tabcolsep}{3pt}\renewcommand{\arraystretch}{1.15}
\begin{longtable}{>{\raggedright}p{0.5cm}>{\raggedright}p{4.0cm}>{\raggedright}p{5.6cm}>{\raggedright}p{2.3cm}>{\raggedright\arraybackslash}p{1.7cm}}
\toprule
N\textsuperscript{o} & Affirmation & Expérience et ce qui est mesuré & Source & Statut \\
\midrule\endfirsthead
\toprule
N\textsuperscript{o} & Affirmation & Expérience et ce qui est mesuré & Source & Statut \\
\midrule\endhead
1 & Les fréquences des neurones \nt{GLUT} du cortex vont d'une fraction à quelques dizaines d'impulsions par seconde, et la plupart du temps les neurones travaillent près de la borne inférieure & Enregistrement par une pipette accolée à la cellule (le choix du neurone ne dépend pas de son activité) dans le cortex auditif du rat éveillé: à tout instant, moins de $5\,\%$ des neurones dépassent $20$ impulsions par seconde; chez une partie des neurones, la fréquence de fond est inférieure à $0.01$ par seconde; la distribution des fréquences est log-normale & \cite{Hromadka2008} & mesuré \\
2 & Capacité d'un canal à impulsions $R_{\max}\approx r\log_2\frac{e}{r\Delta t}$~\eqref{eq:spikeentropy}, presque entièrement dans l'instant des impulsions (proposition~\ref{prop:capacity}) & Entropie d'une chaîne binaire de cases de longueur $\Delta t$. Valeurs du chapitre: $8$ bits par impulsion pour $r=10$ par seconde et $\Delta t=1\ms$, environ $5$ bits pour $\Delta t=10\ms$, moins de $2$ bits par seconde pour un compteur d'impulsions & \cite{MacKay1952} & théorie \\
3 & Les neurones sensoriels transmettent d'un à quelques bits par impulsion, dans les meilleurs cas jusqu'à la moitié de la limite & Reconstruction du stimulus à partir des impulsions de neurones sensoriels dont le stimulus est connu; synthèse des mesures dans le livre & \cite{Rieke1997} & mesuré \\
4 & Le générateur d'impulsions est précis; ce sont les variations rapides de l'entrée qui fixent l'instant précis & Neurones du cortex du rat en tranches: pour un courant répété à fluctuations rapides, les impulsions se répètent à mieux que $1\ms$ près; pour un courant constant, les instants se dispersent & \cite{Mainen1995} & mesuré \\
5 & La synapse n'est pas fiable: plus de la moitié des impulsions n'atteignent pas le destinataire, à cause du caractère aléatoire de la libération & Stimulation minimale des entrées de neurones pyramidaux de l'hippocampe (tranches): plus de la moitié des impulsions ne provoquent pas de réponse. Sont exclues les fluctuations du seuil de stimulation, les défauts de conduction aux ramifications de l'axone et la basse température de l'expérience; il reste la libération probabiliste & \cite{AllenStevens1994} & mesuré \\
6 & Le flux d'impulsions du cortex vivant paraît aléatoire: intervalles dispersés comme dans un flux de Poisson, variance du nombre d'impulsions proche de la moyenne; une partie du \og bruit\fg{} est un message sur les mouvements de l'animal & Enregistrements dans les aires visuelles V1 et MT du singe éveillé: coefficient de variation des intervalles proche de $1$, rapport de la variance du nombre d'impulsions à la moyenne comme pour un processus aléatoire; un modèle qui somme de petites entrées indépendantes donne une dispersion bien plus faible. Dans le cortex visuel de la souris, une part importante de la variabilité se prédit d'après l'enregistrement vidéo des mouvements de l'animal lui-même & \cite{Softky1993,Stringer2019} & mesuré \\
7 & Le code de fréquence est lent: erreur $1/\sqrt{rT}$~\eqref{eq:rateerror}, alors que le cerveau reconnaît une image en $\sim150\ms$ & La formule relève de la statistique de Poisson, déduction du chapitre. Expérience: des sujets décidaient si un animal figurait sur une photographie inconnue montrée pendant $20\ms$; les potentiels cérébraux pour \og oui\fg{} et \og non\fg{} divergent au bout d'environ $150\ms$. Le découpage de ce temps en une dizaine d'étapes de $10$--$20\ms$ est une estimation du chapitre, pas une mesure & \cite{Thorpe1996} & mesuré \\
8 & Le moyennage sur un groupe est limité par la corrélation: l'erreur ne baisse pas de plus d'un facteur $1/\sqrt c$~\eqref{eq:correlated} (proposition~\ref{prop:pooling}) & Paires de neurones de l'aire MT du singe enregistrés simultanément: coefficient de corrélation des nombres d'impulsions de $0.12$ en moyenne, et l'analyse théorique montre que même cette corrélation limite le gain du groupe. Théorie: seule la part de la corrélation qui ressemble au signal impose une limite. Modèle de la position opposée: la fréquence d'un groupe de $50$--$100$ neurones se lit en un seul intervalle entre impulsions, $10$--$50\ms$, et l'instant des impulsions isolées n'est pas utilisé dans le cortex & \cite{Zohary1994,MorenoBote2014,ShadlenNewsome1998} & mesuré \\
9 & Un nombre peut s'écrire dans l'instant de la première impulsion~\eqref{eq:latency}, dans l'ordre des impulsions d'un groupe ou dans la phase d'un rythme local & Rétine de la salamandre: la structure spatiale d'une image brièvement présentée est inscrite dans les délais relatifs des premières impulsions des neurones de sortie, presque indépendamment du contraste. Cellules de lieu de l'hippocampe du rat: en traversant \og leur\fg{} lieu, les impulsions se décalent de plus en plus tôt dans le cycle du rythme thêta ($7$--$12$\,Hz), avec un décalage de $100^\circ$ à $355^\circ$. Dans quelle mesure le cortex utilise l'instant des impulsions reste une question ouverte (ligne~8) & \cite{Gollisch2008,OKeefe1993} & mesuré \\
10 & Le code lu dépend de la constante de temps du destinataire~\eqref{eq:reader} (proposition~\ref{prop:reader}); dans le cortex actif, les neurones sont plus proches de détecteurs de coïncidences & La formule~\eqref{eq:reader} est une déduction du chapitre. Revue d'enregistrements in vivo: dans le cortex actif, la conductance de la membrane est plusieurs fois plus élevée qu'au repos, et la constante de temps d'autant plus courte. Que le cortex change de code précisément ainsi n'est pas montré directement & \cite{Destexhe2003} & hypothèse \\
11 & Une synapse à dépression transmet la variation de fréquence, au fond $\Delta r/r$~\eqref{eq:depression} (fig.~\ref{fig:depression}) & Enregistrements appariés de neurones pyramidaux du cortex: la vitesse de dépression est fixée par la probabilité de libération; avec une dépression lente, la réponse reflète la fréquence, avec une dépression rapide, les coïncidences. Modèle fondé sur la dépression mesurée dans le cortex du rat: des variations relatives égales de fréquence sur des entrées rapides et lentes donnent la même réponse, c'est-à-dire que la synapse transmet $\Delta r/r$. Les nombres $1.7\to2.2$, $6.7$ et $90\ms$ sont des calculs du chapitre & \cite{Tsodyks1997,Abbott1997depression} & mesuré \\
12 & La bouffée est un \og trait\fg: sa probabilité de se perdre, $(1-p)^k$~\eqref{eq:burst}, est faible, et la facilitation la réduit encore & Revue: aux synapses peu fiables, les impulsions isolées se perdent souvent, tandis que les bouffées sont transmises de façon fiable grâce à la facilitation; une bouffée provoque des modifications durables des synapses. Que le cerveau lise la bouffée comme un signe à part est une hypothèse & \cite{Lisman1997,AllenStevens1994} & hypothèse \\
13 & Les neurones \nt{GABA} rapides fixent le rythme et la \og ponctuation\fg; une autre classe encore inhibe les inhibiteurs et ouvre le portillon (proposition~\ref{prop:punctuation}) & Revue des propriétés des classes de neurones \nt{GABA} du cortex. Optogénétique: l'excitation rythmique des neurones \nt{GABA} rapides provoque un rythme gamma, et la réponse à un stimulus dépend de la phase du cycle. Dans le cortex visuel de la souris, les neurones PV s'inhibent mutuellement, les neurones SST inhibent toutes les autres classes, les neurones VIP surtout les SST. L'excitation des neurones VIP chez la souris éveillée lève l'inhibition des neurones \nt{GLUT}; ils sont activés par un signal de renforcement. Le partage des rôles comme règle générale est une hypothèse & \cite{Tremblay2016,Cardin2009,Pfeffer2013,Pi2013} & hypothèse \\
14 & L'énergie limite la fréquence moyenne à environ une impulsion par seconde~\eqref{eq:energy} & Budget énergétique de la substance grise du rongeur d'après des données anatomiques et physiologiques: impulsions et courants synaptiques représentent $47\,\%$ et $34\,\%$ de l'énergie de signalisation; au plus $\sim15\,\%$ des neurones peuvent être actifs simultanément. Pour le cortex humain: au plus $\sim1\,\%$ des neurones peuvent être fortement actifs simultanément. Le nombre de $\sim10^9$ ATP par impulsion et la puissance de $\sim10$\,W pour la signalisation sont des estimations du chapitre fondées sur ces calculs & \cite{Attwell2001,Lennie2003} & théorie \\
15 & Le code est parcimonieux et ajusté pour maximiser le nombre de bits par joule; le cortex échange de la précision contre de l'énergie (proposition~\ref{prop:economy}) & Hypothèse du code économe. Appui: chez des souris soumises à une restriction alimentaire, la conductance des récepteurs AMPA et la dépense synaptique d'ATP baissent dans le cortex visuel ($-29\,\%$), la fréquence des impulsions est conservée, et l'accord sur l'orientation s'élargit de $32\,\%$ & \cite{Barlow1961,Padamsey2022} & hypothèse \\
\bottomrule
\end{longtable}}
